\scanpage{467}
\mspage{379}
% archive sheet number: 466

D'autre part il est immédiat qu'elle est stable par \deleted{\unsure{pe}}petites $\varprojlim$, $N$ l'étant. Donc on conclut \deleted{\unsure{par}} \unsure{encore|en} par 4.8.9 (page 207) que c'est une catégorie paratopos.

\uline{Corollaire} 4.14.9. Soient $M,N$ deux catégories paratopos, et $\underline{\mathrm{Hom}}^{!}(M,N)$ la catégorie des foncteurs $M\longrightarrow N$ qui commutent aux \inserted{petites} $\varprojlim$. On a alors une équivalence de catégories
\[
\begin{aligned}
\underline{\mathrm{Hom}}^{!}(M,N)&\underset{\unsure{-|\leftarrow}}{\approx}\underline{\mathrm{Hom}}_{!}(N,M)^{\circ}\\
\deleted{\text{\unsure{ce qui est|qui associe} \unsure{à tout}}}\qquad f&\longmapsto\text{adjoint à gauche de }f,
\end{aligned}
\]
\margin{\figure{\begin{tikzpicture}[x=.0058cm,y=-.0058cm,line cap=round,line join=round]
\draw[line width=.65pt] (2,30) .. controls (45,1) and (110,-3) .. (149,1) .. controls (202,3) and (217,40) .. (195,83) .. controls (176,123) and (123,156) .. (73,196) .. controls (32,226) and (-2,266) .. (16,286) .. controls (28,302) and (64,304) .. (91,298) .. controls (141,291) and (184,285) .. (229,272);
\draw[line width=.8pt] (2,107) .. controls (4,124) and (4,149) .. (14,167);
\draw[line width=.9pt] (1,176) -- (8,184) -- (11,181);
\draw[line width=.8pt] (244,87) .. controls (247,107) and (246,132) .. (252,149);
\draw[line width=1pt] (254,166) -- (258,176);
\end{tikzpicture}}}%
et par suite la catégorie \underbar{\uuline{opposée}} à $\underline{\mathrm{Hom}}^{!}(M,N)$ est accessible.

Ceci résulte du

\uline{Lemme} 4.14.10 (pour mémoire) \unsure{Soient|Soit} $M$ \deleted{\unsure{et}} \deleted{\unsure{une}} cat.-paratopos, $f:M\longrightarrow N$ un foncteur. Pour que $f$ admette un adjoint à g. (resp à droite) il f. et s. que $f$ commute aux petites $\varprojlim$ \unsure{(|((}resp aux petites $\varinjlim$).
